\chapter{一階述語論理とヒルベルト流証明体系}\label{ch:logic}

本章では本ノートで用いる証明体系 $\mathbf{H}$ を定め, 証明木の書き方と, 証明木を機械的に生成・検査する方法を述べる. 言語と代入は Enderton \cite{enderton} \S 2.1--2.4 に従い, 証明体系は Enderton の論理的公理(\S 2.4)に推論規則として汎化を加えたもの(Monk \cite{monk} 定義 10.23 と同じ型)である. 本章の議論はメタ理論におけるものである.

%%%%%%%%%%%%%%%%%%%%%%%%%%%%%%%%%%%%%%%%%%%%%%%%%%%%%%%%%%%%%%%%%%%%%%
\section{言語と略記}

論理的記号は変数記号 $v_0, v_1, \ldots$ (本文では $x, y, a, \varepsilon, N$ などと書く), 等号 $=$, 否定 $\neg$, 含意 $\to$, 全称量化子 $\forall$ である(Enderton \cite{enderton} \S 2.1). 集合論の言語は $\mathcal{L}_\in = \{\in\}$ であり, 第\ref{ch:definitions}章の方法で記号を順次追加していく. 項と論理式は通常通り帰納的に定義する. 次は略記である(Enderton \S 2.1 と同じ).
\begin{align*}
  (\alpha \land \beta) &:\equiv \neg(\alpha \to \neg\beta), & (\alpha \lor \beta) &:\equiv (\neg\alpha \to \beta), & (\alpha \leftrightarrow \beta) &:\equiv ((\alpha \to \beta) \land (\beta \to \alpha)), \\
  \exists x\,\alpha &:\equiv \neg\forall x\,\neg\alpha, & t \neq u &:\equiv \neg(t = u), & t \notin u &:\equiv \neg(t \in u).
\end{align*}
さらに本ノートでは, 第\ref{ch:definitions}章で導入する関係記号 $<$ について
\[ t \leq u :\equiv (t < u \lor t = u) \]
とし, 有界量化子 $\forall x \in t\,(\alpha) :\equiv \forall x(x \in t \to \alpha)$, $\exists x \in t\,(\alpha) :\equiv \exists x(x \in t \land \alpha)$ を用いる. $\to$ は右結合とし, $\land, \lor$ は $\to$ より強く結合する.

略記は記号列の省略にすぎない. 以下で論理式が等しい, トートロジーである, などと言うときは, すべての略記を展開した原始記号の列について言う.

\begin{dfn}[代入と代入可能性 {\cite[\S 2.4]{enderton}}]
  論理式 $\alpha$ の自由変数 $x$ に項 $t$ を代入したもの $\alpha^x_t$ を, 束縛変数の付け替えをせずに帰納的に定める. $t$ が $\alpha$ の $x$ に\textbf{代入可能}であるとは, $\alpha$ の中で自由に現れる $x$ のどの出現も, $t$ の変数を束縛する量化子 $\forall y$ の作用域の中にないことをいう.
\end{dfn}

%%%%%%%%%%%%%%%%%%%%%%%%%%%%%%%%%%%%%%%%%%%%%%%%%%%%%%%%%%%%%%%%%%%%%%
\section{証明体系 $\mathbf{H}$}

\begin{dfn}[論理的公理 {\cite[\S 2.4]{enderton}}]\label{dfn:hilbert}
  次の形の論理式を\textbf{論理的公理}という. ($x, y$ は変数記号, $\alpha, \beta$ は論理式, $t$ は項.)
  \begin{enumerate}[label=(A\arabic*)]
    \item トートロジー.
    \item $\forall x\,\alpha \to \alpha^x_t$. ただし $t$ は $\alpha$ の $x$ に代入可能.
    \item $\forall x(\alpha \to \beta) \to (\forall x\,\alpha \to \forall x\,\beta)$.
    \item $\alpha \to \forall x\,\alpha$. ただし $x$ は $\alpha$ に自由に現れない.
    \item $x = x$.
    \item $x = y \to (\alpha \to \alpha')$. ただし $\alpha$ は原子論理式で, $\alpha'$ は $\alpha$ の $x$ のいくつか(0個でもよい)を $y$ に置き換えたもの.
  \end{enumerate}
  ここで\textbf{トートロジー}とは, 原子論理式と $\forall x\,\alpha$ の形の論理式(\textbf{素論理式})を命題変数とみなしたとき, $\neg, \to$ の真理値表のもとで常に真になる論理式である.
\end{dfn}

\begin{dfn}[推論規則]
  \begin{enumerate}[label=(R\arabic*)]
    \item (MP) $\alpha$ と $\alpha \to \beta$ から $\beta$ を導く.
    \item (Gen) $\alpha$ から $\forall x\,\alpha$ を導く.
  \end{enumerate}
\end{dfn}

\begin{dfn}[証明・定理]
  $T$ を閉論理式の集合(理論)とする. 論理式の有限列 $\alpha_1, \ldots, \alpha_n$ が $T$ からの $\alpha_n$ の\textbf{証明}であるとは, 各 $\alpha_i$ が論理的公理であるか, $T$ の元であるか, それより前の論理式から MP または Gen で得られることをいう. 存在するとき $T \vdash_{\mathbf{H}} \alpha_n$ と書き, $\alpha_n$ を $T$ の\textbf{定理}という.
\end{dfn}

\begin{rem}
  $T$ が閉論理式のみからなるので, Gen をどこで用いても健全である. Enderton の体系では推論規則は MP のみで, 代わりに「論理的公理の汎化もまた論理的公理」とする. 両者は同じ定理を与える: $\mathbf{H}$ の Gen は Enderton の汎化定理(定理 24B. $T$ の元は閉論理式なので変数条件は自動的にみたされる)により Enderton の体系で許容され, 逆に Enderton の論理的公理(公理の汎化)は $\mathbf{H}$ で公理に Gen を施して得られる. 従って健全性定理・完全性定理(Enderton \S 2.5)はそのまま $\mathbf{H}$ に適用できる. 推論規則として汎化をもつ同型の体系は Monk \cite{monk} 定義 10.23 にもある.
\end{rem}

\begin{dfn}[証明木]\label{dfn:prooftree}
  $T$ からの $\alpha$ の\textbf{証明木}とは, 各ノードが論理式でラベル付けされた有限の根付き木で, 根のラベルが $\alpha$ であり, 各ノードが次のいずれかであるものをいう.
  \begin{enumerate}[label=(\arabic*)]
    \item 葉であって, ラベルが論理的公理, $T$ の元, または既に証明された定理である.
    \item 2つの子をもち, 左の子のラベルが $\beta$, 右の子のラベルが $\beta \to \gamma$, 自身のラベルが $\gamma$ である (MP).
    \item 1つの子をもち, 子のラベルが $\beta$, 自身のラベルが $\forall x\,\beta$ である (Gen).
  \end{enumerate}
\end{dfn}

証明木を子が親より先に来るように並べれば証明が得られ, 逆も成り立つ. 葉に既証の定理を置くことは, その定理の証明木をそこに接ぎ木することの省略であり, 定理を証明した順序が循環しない限り許される.

\section{証明木の読み方}\label{sec:reading}

本ノートの「完全展開した証明木」はすべて, 第\ref{sec:generator}節の生成器が出力し, 機械的に検査したものである. 次の約束で印刷する.
\begin{itemize}
  \item 推論規則の名前は書かない. \textbf{二分岐のノードは MP}(左が $\beta$, 右が $\beta \to \gamma$), \textbf{一分岐のノードは Gen} である.
  \item 葉の右に小さく根拠を書く: \textsf{A1}--\textsf{A6} は論理的公理, \textsf{Ext}, \textsf{Sep} は $\ZFC$ の公理, \textsf{Def$\cdots$} は定義公理, [O1] や 定理2.3 などは既証の定理, (R$k$), [N$k$], [S$k$] は付録の基礎定理である.
  \item 木が大きいときは部分木を切り出して $\Pi_1, \Pi_2, \ldots$ と名前を付け, 先に印刷する. 使う場所には $\overset{\Pi_k}{\vdots}$ とその根の論理式を書く.
  \item 長い論理式は $\varphi_1, \varphi_2, \ldots$ と略記し, 木の後に「略記」として一覧を付ける. 略記の定義の中に他の略記が現れることがある. これらは記号列の省略であり, 証明の内容には影響しない.
  \item (A1) の葉は $\beta_1 \to (\beta_2 \to \cdots \to (\beta_n \to \gamma))$ の形のトートロジーで, その下に MP が $n$ 回続き, $\beta_1, \ldots, \beta_n$ から命題論理だけで $\gamma$ を導く(Enderton の Rule T, 補題 24C).
\end{itemize}
例えば, 仮定 $h_1, \ldots, h_k$ のもとで $\varphi$ を示す議論は, 木の中では $h_1 \to (h_2 \to \cdots (h_k \to \varphi))$ という一つの論理式として現れる(下の \ref{sec:generator} 節).

%%%%%%%%%%%%%%%%%%%%%%%%%%%%%%%%%%%%%%%%%%%%%%%%%%%%%%%%%%%%%%%%%%%%%%
\section{証明木の生成と検査}\label{sec:generator}

完全展開した証明木は非常に大きくなる(例えば三角不等式で 155 ノード, アルキメデスの公理で 544 ノード)ので, 人手で書くと誤りが避けられない. そこで次のプログラムを作成した(\texttt{LaTeX/calculus/prover/}, Node.js).

\begin{description}
  \item[\texttt{core.js}] 項・論理式・代入・代入可能性を Enderton の定義通りに実装し, 各ノードが (A1)--(A6), 公理, 定義公理, 既証定理, MP, Gen のいずれかであることを検査する. (A1) の判定は, 素論理式を命題変数とみなした論理式の否定が充足不能であることを, Tseitin 変換と DPLL 法で確かめて行う. (A2) は $t$ を探して $\alpha^x_t$ との一致と代入可能性を, (A4) は自由変数の条件を, (A6) は置き換えの形を検査する.
  \item[\texttt{tactics.js}] 証明木を組み立てる手続き. 以下の「構成法」はメタ理論の手続きであり, それが作る木は公理・MP・Gen だけからなる.
  \item[\texttt{p\_*.js}] 各定理の証明の記述. 最後に \texttt{verify} が木全体をもう一度検査してから \LaTeX{} を出力する.
\end{description}

主な構成法は次の通りである(いずれも Enderton \S 2.4 のメタ定理の証明を, そのまま木の構成として実行したもの).
\begin{enumerate}[label=(\arabic*)]
  \item \textbf{文脈}: 仮定 $h_1, \ldots, h_k$ のもとでの行 $\varphi$ を, 論理式 $h_1 \to \cdots \to h_k \to \varphi$ の証明木として持つ. 仮定を一つ解消すること(演繹定理)は, 同じ論理式の読み方を変えるだけなので木を変えない.
  \item \textbf{命題論理}: 行 $\beta_1, \ldots, \beta_n$ から $\gamma$ がトートロジー的に従うとき, (A1) の葉 $\beta_1 \to \cdots \to \beta_n \to \gamma$ と MP $n$ 回.
  \item \textbf{全称例化}: (A2) の葉と MP. 代入可能でないときは, 先頭の束縛変数を新しい変数で例化してから Gen し直して付け替える(Enderton 定理 24I の具体的な実行).
  \item \textbf{仮定のもとでの汎化}: $H \to \varphi$ ($x$ は $H$ に自由に現れない)から Gen で $\forall x(H \to \varphi)$, (A3) で $\forall x H \to \forall x\varphi$, (A4) で $H \to \forall x H$ を得て, 命題論理で $H \to \forall x\varphi$.
  \item \textbf{存在導入}: $\varphi^x_t$ と (A2) の葉 $\forall x\neg\varphi \to \neg\varphi^x_t$ から命題論理で $\neg\forall x\neg\varphi$, すなわち $\exists x\varphi$.
  \item \textbf{存在除去}: $\exists x\varphi$ と, 仮定 $\varphi$ のもとでの $\psi$ ($x$ は $\psi$ と他の仮定に自由に現れない)から, $(H \land \neg\psi) \to \neg\varphi$ に Gen, (A3), (A4) を施し, 命題論理で $\psi$.
  \item \textbf{等号による書き換え}: 原子論理式 $\beta(z)$ ($z$ は新しい変数)について (A6) の葉 $z = w \to (\beta(z) \to \beta(w))$ に Gen を2回, (A2) で $z := t$, $w := u$ と例化して $t = u \to (\beta(t) \to \beta(u))$ を得る. 対称律 $\forall x \forall y(x = y \to y = x)$ は (A6) $x = y \to (x = x \to y = x)$ と (A5) から, 推移律も (A6) から作る.
  \item \textbf{束縛変数の付け替え}: 素論理式 $\forall u\,\alpha$ と $\forall w\,\alpha^u_w$ が同時に現れて命題論理で結べないときは, (A2), Gen, (A3), (A4) で $\forall u\,\alpha \to \forall w\,\alpha^u_w$ を作って補う.
\end{enumerate}

通常の証明の文章では, これらの構成法を順に \textbf{文脈}, \textbf{TC} (トートロジー的帰結), \textbf{$\forall$E} (全称例化), \textbf{Gen}, \textbf{$\exists$I}, \textbf{$\exists$E}, \textbf{Eq} (等号による書き換え), \textbf{付け替え} と呼ぶことがある. これらはすべて木を作る手続きであって, 木の中に「導出規則」として残るものはない. 生成された木は印刷される前に, すべてのノードについて定義\ref{dfn:prooftree}の条件が再検査される.

\begin{rem}[図式]
  数学的帰納法のような\textbf{定理図式}は, 任意の論理式 $\varphi(n)$ の代わりに1項関係記号 $\Phi$ を用いて証明木を作る. $\Phi(t)$ を $\varphi(t)$ に一斉に置き換えると(必要なら束縛変数を付け替えて), 木の各ノードは同じ種類の公理・推論のままであり, 求める定理の証明木が得られる(Enderton \S 2.4 の議論と同様).
\end{rem}
